\section{The Wake}

The knit moves tones on existing docks. A second law, the wake, governs the open edge: new docks are born there, quiet,
at the resting value of peace, pushing the boundary outward beat by beat. The wake is the law of space, as the knit is
the law of state.

The wake costs no added energy. It is simply what a mesh growing at its edge looks like from within, and it is the
source of cosmic expansion. It also caps the curvature, so the tone never piles up without bound, and there is no initial
singularity.

Now the apparent paradox and its resolution. The local law is perfectly reversible and runs the same forward and
backward, so how can there be a one-way arrow of time. The answer is that the arrow is not in the local law, it is in the
global growth. The knit, microscopically, has no preferred direction, but the mesh as a whole only ever gets larger,
because the wake keeps adding docks at the edge and never removes any, and because distinctions, once made, only
accumulate, since to erase a distinction is itself a distinction, the difference between before and after.

So phase-space volume is conserved by the reversible law while the accessible volume grows, and the result is a clear
macroscopic before and after on a microscopically time-symmetric law. Time has a definite shape, a growing block: the
past is the settled interior, fixed and done, the future is open at the wake where docks are not yet laid down, and the
present is the growing edge. This is the integer order from the start of the construction, made dynamical.
